\chapter{Puissances}

\noindent\textit{Les puissances permettent d'écrire simplement des produits répétés, mais
aussi des nombres extrêmement grands (distances astronomiques) ou extrêmement petits (tailles
microscopiques). Après avoir revu les règles de calcul sur les puissances, tu apprendras à
utiliser la \textbf{notation scientifique} et à estimer un \textbf{ordre de grandeur} — des
outils indispensables en sciences.}
\vspace{8pt}

\connecteBanner

\section{Puissances d'exposant entier positif}

\begin{definition}[Puissance d'exposant positif]
Pour $a$ un nombre et $n$ un entier naturel non nul, $a^n$ (lire « $a$ exposant $n$ » ou
« $a$ puissance $n$ ») est le produit de $n$ facteurs égaux à $a$ :
\[
a^{n}=\underbrace{a\times a\times\cdots\times a}_{n\text{ facteurs}}.
\]
On convient de plus que $a^{1}=a$ et, pour $a\ne0$, $a^{0}=1$.
\end{definition}

\begin{exemple}
$2^{5}=2\times2\times2\times2\times2=32$ ;\quad $(-3)^{4}=81$ ;\quad $(-2)^{5}=-32$ ;\quad
$7^{0}=1$.
\end{exemple}

\begin{remarque}
Le signe d'une puissance d'un nombre négatif dépend de la \textbf{parité} de l'exposant :
si $n$ est \textbf{pair}, $(-a)^{n}>0$ ; si $n$ est \textbf{impair}, $(-a)^{n}<0$ (pour
$a>0$).
\end{remarque}

\begin{illus}
\begin{tikzpicture}
\draw[fill=onbsoft,draw=onbo,line width=0.9pt] (0,0) rectangle (0.5,0.5);
\node[below] at (0.25,-0.15) {\small $1^{2}=1$};
\draw[fill=onbsoft,draw=onbo,line width=0.9pt] (1.2,0) rectangle (2.2,1);
\node[below] at (1.7,-0.15) {\small $2^{2}=4$};
\draw[fill=onbsoft,draw=onbo,line width=0.9pt] (3,0) rectangle (4.5,1.5);
\node[below] at (3.75,-0.15) {\small $3^{2}=9$};
\draw[fill=onbsoft,draw=onbo,line width=0.9pt] (5.2,0) rectangle (7.2,2);
\node[below] at (6.2,-0.15) {\small $4^{2}=16$};
\end{tikzpicture}
\fig{Figure — L'aire d'un carré de côté $n$ vaut $n^2$ : elle croît très vite avec $n$.}
\end{illus}

\section{Puissances d'exposant négatif et nul}

\begin{definition}[Puissance d'exposant négatif]
Pour $a\ne0$ et $n$ un entier naturel non nul, on pose
\[
a^{-n}=\frac{1}{a^{n}}.
\]
\end{definition}

\begin{exemple}
$2^{-3}=\dfrac1{2^{3}}=\dfrac18$ ;\quad $5^{-2}=\dfrac1{25}$ ;\quad
$(-3)^{-2}=\dfrac1{(-3)^{2}}=\dfrac19$ ;\quad $10^{-4}=\dfrac1{10\,000}$.
\end{exemple}

\begin{remarque}
$0$ n'a pas d'exposant négatif ($\dfrac1{0^{n}}$ n'existe pas) : $0^{-n}$ n'est jamais défini.
En revanche $0^{n}=0$ pour tout entier $n\ge1$, et $a^{0}=1$ pour tout $a\ne0$ (y compris si
$a$ est négatif).
\end{remarque}

\section{Règles de calcul sur les puissances}

\begin{propriete}[Produit et quotient de puissances de même base]
Pour $a\ne0$ et $m,n$ entiers relatifs :
\[
a^{m}\times a^{n}=a^{m+n}, \qquad \frac{a^{m}}{a^{n}}=a^{m-n}.
\]
\end{propriete}

\begin{propriete}[Puissance d'une puissance, d'un produit, d'un quotient]
Pour $a,b\ne0$ et $m,n$ entiers relatifs :
\[
\left(a^{m}\right)^{n}=a^{m\times n}, \qquad (a\times b)^{n}=a^{n}\times b^{n}, \qquad
\left(\frac ab\right)^{n}=\frac{a^{n}}{b^{n}}.
\]
\end{propriete}

\begin{illus}
\begin{tikzpicture}[>=Stealth,every node/.style={font=\small}]
\node[draw=onbo,fill=onbsoft,rounded corners=3pt,minimum width=2.5cm,minimum height=0.7cm] (a1) {$a^m\times a^n$};
\node[draw=onbblue,fill=white,rounded corners=3pt,minimum width=2.2cm,minimum height=0.7cm,right=14mm of a1] (a2) {$a^{m+n}$};
\draw[->,onbmuted,line width=1pt] (a1) -- (a2);
\node[draw=onbo,fill=onbsoft,rounded corners=3pt,minimum width=2.5cm,minimum height=0.7cm,below=8mm of a1] (b1) {$a^m\div a^n$};
\node[draw=onbblue,fill=white,rounded corners=3pt,minimum width=2.2cm,minimum height=0.7cm,right=14mm of b1] (b2) {$a^{m-n}$};
\draw[->,onbmuted,line width=1pt] (b1) -- (b2);
\node[draw=onbo,fill=onbsoft,rounded corners=3pt,minimum width=2.5cm,minimum height=0.7cm,below=8mm of b1] (c1) {$(a^m)^n$};
\node[draw=onbblue,fill=white,rounded corners=3pt,minimum width=2.2cm,minimum height=0.7cm,right=14mm of c1] (c2) {$a^{m\times n}$};
\draw[->,onbmuted,line width=1pt] (c1) -- (c2);
\end{tikzpicture}
\fig{Figure — Les trois règles de base sur les exposants.}
\end{illus}

\begin{exemple}
$3^{2}\times3^{4}=3^{6}=729$ ;\quad $5^{7}\div5^{3}=5^{4}=625$ ;\quad
$\left(2^{3}\right)^{4}=2^{12}=4\,096$.
\end{exemple}

\begin{illus}
\begin{tikzpicture}[scale=0.4]
\draw[onbline,line width=0.4pt] (0,0) grid (6,6);
\draw[onbo,line width=1.4pt] (0,0) rectangle (6,6);
\draw[onbo,line width=1.4pt] (3,0) -- (3,6);
\draw[onbo,line width=1.4pt] (0,3) -- (6,3);
\end{tikzpicture}
\fig{Figure — $(2\times3)^{2}=36$ cases, regroupées en $4$ blocs de $9$ :
$2^{2}\times3^{2}=36$.}
\end{illus}

\begin{aretenir}
On \textbf{n'applique} ces règles que si les puissances ont la \textbf{même base} (pour la
somme/différence des exposants) : $2^{3}\times5^{2}$ ne se simplifie \textbf{pas} en une seule
puissance. Ne jamais confondre $a^{m}\times a^{n}=a^{m+n}$ (on \textbf{additionne}) avec
$\left(a^{m}\right)^{n}=a^{m\times n}$ (on \textbf{multiplie}).
\end{aretenir}

\section{Puissances de 10, notation scientifique}

\begin{propriete}[Puissances de 10]
Pour tout entier naturel $n$, $10^{n}$ s'écrit $1$ suivi de $n$ zéros, et
$10^{-n}=\dfrac1{10^{n}}$ s'écrit avec $n$ zéros après la virgule (dont le dernier est le
chiffre $1$). Multiplier par $10^{n}$ revient à déplacer la virgule de $n$ rangs vers la
droite ; par $10^{-n}$, de $n$ rangs vers la gauche.
\end{propriete}

\begin{definition}[Notation scientifique]
La \textbf{notation scientifique} d'un nombre décimal non nul est son écriture sous la forme
\[
a\times10^{n}, \qquad 1\le|a|<10,\ n\in\Z.
\]
\end{definition}

\begin{exemple}
$1\,250\,000=1{,}25\times10^{6}$ ;\quad $45\,600\,000\,000=4{,}56\times10^{10}$ ;\quad
$0{,}000034=3{,}4\times10^{-5}$ ;\quad $0{,}0000000078=7{,}8\times10^{-9}$.
\end{exemple}

\begin{illus}
\begin{tikzpicture}[>=Stealth]
\draw[onbline,line width=1pt] (0,0) -- (11,0);
\foreach \x/\e in {0/-7,2.75/-3,5.5/0,8.25/7,11/11}
  \draw[onbmuted] (\x,0.08) -- (\x,-0.08) node[below] {\small $10^{\e}$};
\filldraw[onbgreen] (0,0) circle(2pt) node[above,yshift=1pt,text width=1.7cm,align=center] {\tiny virus};
\filldraw[onbo] (2.75,0) circle(2pt) node[above,yshift=1pt,text width=1.9cm,align=center] {\tiny grain de sable};
\filldraw[onbblue] (5.5,0) circle(2pt) node[above,yshift=1pt,text width=1.9cm,align=center] {\tiny être humain};
\filldraw[onbo2] (8.25,0) circle(2pt) node[above,yshift=1pt,text width=2.1cm,align=center] {\tiny population du Cameroun};
\filldraw[onbink] (11,0) circle(2pt) node[above,yshift=1pt,text width=1.9cm,align=center] {\tiny distance Terre-Soleil};
\end{tikzpicture}
\fig{Figure — Échelle des ordres de grandeur, de l'infiniment petit à l'infiniment
grand.}
\end{illus}

\section{Ordre de grandeur}

\begin{definition}[Ordre de grandeur]
L'\textbf{ordre de grandeur} d'un nombre est la puissance de $10$ la plus proche de ce
nombre. On l'obtient en général à partir de la notation scientifique $a\times10^{n}$ : si
$a<5$, l'ordre de grandeur est $10^{n}$ ; si $a\ge5$, c'est $10^{n+1}$.
\end{definition}

\begin{exemple}
$2\,937{,}6=2{,}9376\times10^{3}$ : comme $2{,}9376<5$, l'ordre de grandeur est $10^{3}$.
$3{,}7\times10^{5}\times2{,}9\times10^{-8}=10{,}73\times10^{-3}=1{,}073\times10^{-2}$ : ordre
de grandeur $10^{-2}$.
\end{exemple}

\begin{aretenir}
L'ordre de grandeur sert à \textbf{vérifier rapidement} la vraisemblance d'un résultat de
calcul, sans calculatrice, en estimant chaque facteur par la puissance de $10$ la plus proche.
\end{aretenir}

% ═══════════════════════════════════════════════════════════════
\section{L'essentiel à retenir}

\begin{synthese}
\textbf{Puissance positive.} $a^{n}=a\times\cdots\times a$ ($n$ facteurs) ; $a^{0}=1$ (pour
$a\ne0$).\\[4pt]
\textbf{Puissance négative.} $a^{-n}=\dfrac1{a^{n}}$ (pour $a\ne0$).\\[4pt]
\textbf{Règles (même base).} $a^{m}\times a^{n}=a^{m+n}$ ;\
$\dfrac{a^{m}}{a^{n}}=a^{m-n}$ ;\ $(a^{m})^{n}=a^{mn}$.\\[4pt]
\textbf{Puissance d'un produit/quotient.} $(ab)^{n}=a^{n}b^{n}$ ;\
$\left(\dfrac ab\right)^{n}=\dfrac{a^{n}}{b^{n}}$.\\[4pt]
\textbf{Notation scientifique.} $a\times10^{n}$ avec $1\le|a|<10$. Multiplier par $10^{n}$
déplace la virgule de $n$ rangs (droite si $n>0$, gauche si $n<0$).\\[4pt]
\textbf{Ordre de grandeur.} La puissance de $10$ la plus proche du nombre (arrondir le
coefficient $a$ de la notation scientifique).\\[4pt]
\textbf{Piège fréquent.} $a^{m}\times a^{n}=a^{m+n}$ (on additionne) $\ne$
$(a^{m})^{n}=a^{m\times n}$ (on multiplie). Les règles de base ne s'appliquent qu'entre
puissances de \textbf{même base}.
\end{synthese}

% ═══════════════════════════════════════════════════════════════
\section{Méthodes \& savoir-faire}

\methode{Appliquer les règles de calcul sur des puissances de même base}
Repérer si l'opération est un produit, un quotient ou une puissance de puissance, puis
additionner, soustraire ou multiplier les exposants en conséquence.
\begin{exemple}
$\dfrac{2^{3}\times2^{6}}{2^{4}}=\dfrac{2^{9}}{2^{4}}=2^{5}=32$.
\end{exemple}

\methode{Calculer une puissance d'exposant négatif}
Écrire $a^{-n}=\dfrac1{a^{n}}$, calculer $a^{n}$, puis prendre l'inverse.
\begin{exemple}
$4^{-3}=\dfrac1{4^{3}}=\dfrac1{64}$.
\end{exemple}

\methode{Écrire un nombre en notation scientifique}
Déplacer la virgule pour obtenir un nombre $a$ tel que $1\le|a|<10$ ; le nombre de rangs
déplacés (positif si on l'a déplacée vers la gauche, négatif vers la droite) donne
l'exposant $n$.
\begin{exemple}
$523\,000=5{,}23\times10^{5}$ (virgule déplacée de 5 rangs vers la gauche) ;
$0{,}0091=9{,}1\times10^{-3}$ (virgule déplacée de 3 rangs vers la droite).
\end{exemple}

\methode{Passer de la notation scientifique à l'écriture décimale}
Déplacer la virgule de $|n|$ rangs : vers la droite si $n>0$, vers la gauche si $n<0$ (en
complétant par des zéros si nécessaire).
\begin{exemple}
$7{,}4\times10^{4}=74\,000$ ;\quad $6{,}05\times10^{-3}=0{,}00605$.
\end{exemple}

\methode{Calculer un produit ou un quotient en notation scientifique}
Regrouper les coefficients entre eux et les puissances de $10$ entre elles, puis remettre le
résultat sous forme scientifique si nécessaire.
\begin{exemple}
$(3\times10^{4})\times(2\times10^{-7})=6\times10^{-3}$ ;\quad
$(8\times10^{6})\div(4\times10^{9})=2\times10^{-3}$.
\end{exemple}

\methode{Encadrer / estimer un ordre de grandeur}
Remplacer chaque nombre par sa notation scientifique arrondie, puis calculer avec ces valeurs
approchées.
\begin{exemple}
$598\times0{,}0021\approx6\times10^{2}\times2\times10^{-3}=12\times10^{-1}=1{,}2$ : ordre de
grandeur $1$ (valeur exacte : $1{,}2558$).
\end{exemple}

\methode{Simplifier une expression combinant plusieurs puissances de même base}
Réécrire chaque facteur avec la même base, puis appliquer les règles dans l'ordre (produits
et quotients d'abord, puissance de puissance ensuite si besoin).
\begin{exemple}
$\dfrac{5^{3}\times5^{-6}}{5^{-4}}=5^{3-6-(-4)}=5^{1}=5$.
\end{exemple}

% ═══════════════════════════════════════════════════════════════
\section{Exercices d'application}
\begin{multicols}{2}

\rubrique{Calculer avec les règles de puissances}
\exo{}\diff{1} Écrire sous forme d'une seule puissance puis calculer : $3^{2}\times3^{4}$ ;\
$(-2)^{3}\times(-2)^{5}$ ;\ $5^{7}\div5^{3}$ ;\ $\left(2^{3}\right)^{4}$.

\exo{}\diff{1} Calculer : $7^{0}$ ;\ $(-5)^{1}$ ;\ $(-1)^{7}$ ;\ $(-1)^{10}$.

\exo{}\diff{2} Calculer : $2^{-3}$ ;\ $5^{-2}$ ;\ $(-3)^{-2}$ ;\ $10^{-4}$.

\exo{}\diff{2} Calculer : $3^{4}\times3^{-2}$ ;\ $2^{-3}\times2^{5}$ ;\ $\left(5^{2}\right)^{-1}$ ;\ $\left(2^{-1}\right)^{3}$.

\rubrique{Notation scientifique}
\exo{}\diff{1} Écrire en notation scientifique : $1\,250\,000$ ;\ $45\,600\,000\,000$.

\exo{}\diff{1} Écrire en notation scientifique : $0{,}000034$ ;\ $0{,}0000000078$.

\exo{}\diff{2} Donner l'écriture décimale de $6{,}02\times10^{4}$ et de $3{,}1\times10^{-3}$.

\exo{}\diff{2} Calculer, en donnant le résultat en notation scientifique :
$(2{,}4\times10^{5})\times(3\times10^{-2})$ et $(6\times10^{8})\div(1{,}5\times10^{3})$.

\rubrique{Ordre de grandeur}
\exo{}\diff{2} Donner un ordre de grandeur de $3{,}7\times10^{5}\times2{,}9\times10^{-8}$.

\exo{}\diff{2} Donner un ordre de grandeur de $612\times4{,}8$, sans calculatrice, puis
vérifier avec une calculatrice.

\exo{}\diff{2} Encadrer $2^{10}$ entre deux puissances de $10$ consécutives ; en déduire un
ordre de grandeur de $2^{10}$.

\exo{}\diff{3} Un fichier informatique pèse $2^{20}$ octets. Sachant que $2^{10}\approx10^{3}$,
donner une valeur approchée de $2^{20}$ en notation scientifique, puis calculer sa valeur
exacte.

\rubrique{Problèmes concrets (astronomie, biologie)}
\exo{}\diff{2} La distance moyenne Terre-Soleil est $149\,600\,000$ km. Écrire cette distance
en notation scientifique, en km puis en m.

\exo{}\diff{2} Un globule rouge mesure environ $0{,}000007$ m de diamètre. Écrire cette taille
en notation scientifique, puis calculer le nombre de globules rouges qu'il faudrait aligner
pour obtenir $1$ cm.

\exo{}\diff{2} Le Cameroun compte environ $2{,}7\times10^{7}$ habitants, et la population
mondiale environ $8\times10^{9}$ habitants. Donner un ordre de grandeur de la fraction que
représente la population camerounaise dans la population mondiale.

\exo{}\diff{1} La masse de la Terre est environ $5{,}97\times10^{24}$ kg. Exprimer cette masse
en tonnes ($1$ tonne $=10^{3}$ kg), en notation scientifique.

% ═══════════════════════════════════════════════════════════════
\end{multicols}

\section{Exercices d'entraînement}
\begin{multicols}{2}

\rubrique{Calculs avec les règles}
\exo{}\diff{1} Calculer, sous forme d'une seule puissance puis d'un nombre : $3^{3}\times3^{5}$.

\exo{}\diff{2} Calculer : $(-4)^{2}\times(-4)^{3}$.

\exo{}\diff{1} Calculer : $2^{8}\div2^{5}$.

\exo{}\diff{2} Calculer : $\left(5^{2}\right)^{3}$.

\exo{}\diff{2} Calculer : $\dfrac{7^{5}\times7^{2}}{7^{4}}$.

\rubrique{Puissances d'exposant négatif}
\exo{}\diff{1} Calculer : $4^{-2}$ ;\ $3^{-3}$.

\exo{}\diff{2} Calculer : $(-2)^{-3}$ ;\ $10^{-5}$.

\exo{}\diff{2} Calculer : $2^{-2}\times2^{6}$.

\exo{}\diff{2} Calculer : $\left(3^{-1}\right)^{2}$ et $\left(7^{3}\right)^{-1}$.

\rubrique{Notation scientifique}
\exo{}\diff{1} Écrire en notation scientifique : $860\,000$ ;\ $7\,100\,000\,000$.

\exo{}\diff{2} Écrire en notation scientifique : $0{,}0000523$ ;\ $0{,}00000000042$.

\exo{}\diff{1} Donner l'écriture décimale de $4{,}5\times10^{5}$ et de $2{,}3\times10^{-4}$.

\exo{}\diff{2} Donner l'écriture décimale de $9{,}02\times10^{6}$ et de $5\times10^{-7}$.

\exo{}\diff{2} Calculer, en notation scientifique : $(3\times10^{4})\times(2\times10^{3})$ et
$(8\times10^{6})\div(4\times10^{2})$.

\rubrique{Ordre de grandeur}
\exo{}\diff{2} Donner un ordre de grandeur de $5\times10^{-3}\times6\times10^{5}$.

\exo{}\diff{2} Donner un ordre de grandeur de $9\times10^{-2}\div3\times10^{-5}$.

\exo{}\diff{2} Donner un ordre de grandeur de $4\,987\times0{,}0031$.

\exo{}\diff{3} Encadrer $3^{7}$ entre deux puissances de $10$ consécutives.

\rubrique{Problèmes concrets}
\exo{}\diff{2} La distance Terre-Lune est environ $384\,000$ km. Écrire cette distance en
notation scientifique, en km puis en m.

\exo{}\diff{2} Un cheveu humain a un diamètre d'environ $0{,}00007$ m. L'écrire en notation
scientifique et donner un ordre de grandeur en $\mu$m ($1\,\mu\text{m}=10^{-6}$ m).

\exo{}\diff{2} Une bactérie mesure environ $2\times10^{-6}$ m. Combien de bactéries alignées
faudrait-il pour atteindre $1$ mm ?

\exo{}\diff{2} La superficie du Cameroun est d'environ $475\,000\text{ km}^{2}$ et celle de
l'Afrique d'environ $30\,000\,000\text{ km}^{2}$. Donner un ordre de grandeur de la fraction
que représente le Cameroun sur le continent africain.

\exo{}\diff{3} Un ordinateur effectue environ $3\times10^{9}$ opérations par seconde. Combien
d'opérations effectue-t-il en une heure ? Donner le résultat en notation scientifique.

% ═══════════════════════════════════════════════════════════════
\end{multicols}

\section{Sujets type BEPC}

\sujetbac[puissances et lumière du Soleil]
\begin{enumerate}[label=\arabic*)]
\item Calculer, et donner le résultat sous forme d'une puissance de $2$ :
$D=\dfrac{2^{3}\times2^{-5}}{2^{-4}}$. En déduire sa valeur numérique.
\item Écrire en notation scientifique la distance moyenne Terre-Soleil,
$149\,600\,000$ km.
\item Sachant que la lumière parcourt environ $3\times10^{5}$ km chaque seconde, calculer le
temps (en secondes, arrondi au dixième) mis par la lumière du Soleil pour atteindre la Terre.
Convertir ce temps en minutes et secondes (arrondir les secondes à l'unité).
\item Calculer $E=\left(4\times10^{3}\right)^{2}$ et donner le résultat en notation
scientifique.
\end{enumerate}

\sujetbac[puissances de 2 et stockage numérique — cybercafé de Yaoundé]
Un cybercafé de Yaoundé équipe ses ordinateurs de disques durs de capacité $2^{30}$ octets.
\begin{enumerate}[label=\arabic*)]
\item Sachant que $2^{10}=1\,024$, calculer les valeurs exactes de $2^{20}$ et de $2^{30}$.
\item Chaque photo enregistrée pèse $2^{22}$ octets. Exprimer, sous forme d'une puissance de
$2$, le nombre maximal de photos que peut contenir un disque dur de $2^{30}$ octets ; en
déduire sa valeur exacte.
\item Écrire $2^{30}$ en notation scientifique (arrondir le coefficient au centième).
\item Donner un ordre de grandeur, en octets, de la capacité de ce disque dur.
\end{enumerate}

\sujetbac[puissances de 10 et population du Cameroun]
\begin{enumerate}[label=\arabic*)]
\item Calculer, sous forme d'une seule puissance de $10$ :
$F=\dfrac{10^{7}\times10^{-2}}{10^{-3}}$.
\item Le Cameroun compte environ $2{,}7\times10^{7}$ habitants et la population mondiale
environ $8\times10^{9}$ habitants.
\begin{enumerate}[label=\alph*)]
\item Calculer, en notation scientifique (coefficient arrondi au dixième), la fraction que
représente la population du Cameroun dans la population mondiale.
\item En déduire le pourcentage, arrondi au dixième, que représente la population
camerounaise dans la population mondiale.
\end{enumerate}
\item Le Cameroun a une superficie d'environ $475\,000\text{ km}^{2}$. Écrire cette
superficie en notation scientifique.
\end{enumerate}

% ═══════════════════════════════════════════════════════════════
\section{Corrigés}
\begin{multicols}{2}

\corrige{de l'exercice 1}
$3^{2}\times3^{4}=3^{6}=729$ ;\ $(-2)^{3}\times(-2)^{5}=(-2)^{8}=256$ ;\
$5^{7}\div5^{3}=5^{4}=625$ ;\ $\left(2^{3}\right)^{4}=2^{12}=4\,096$.

\corrige{de l'exercice 2}
$7^{0}=1$ ;\ $(-5)^{1}=-5$ ;\ $(-1)^{7}=-1$ ;\ $(-1)^{10}=1$.

\corrige{de l'exercice 3}
$2^{-3}=\dfrac18$ ;\ $5^{-2}=\dfrac1{25}$ ;\ $(-3)^{-2}=\dfrac19$ ;\ $10^{-4}=\dfrac1{10\,000}$.

\corrige{de l'exercice 4}
$3^{4}\times3^{-2}=3^{2}=9$ ;\ $2^{-3}\times2^{5}=2^{2}=4$ ;\
$\left(5^{2}\right)^{-1}=5^{-2}=\dfrac1{25}$ ;\ $\left(2^{-1}\right)^{3}=2^{-3}=\dfrac18$.

\corrige{de l'exercice 5}
$1\,250\,000=1{,}25\times10^{6}$ ;\ $45\,600\,000\,000=4{,}56\times10^{10}$.

\corrige{de l'exercice 6}
$0{,}000034=3{,}4\times10^{-5}$ ;\ $0{,}0000000078=7{,}8\times10^{-9}$.

\corrige{de l'exercice 7}
$6{,}02\times10^{4}=60\,200$ ;\ $3{,}1\times10^{-3}=0{,}0031$.

\corrige{de l'exercice 8}
$(2{,}4\times10^{5})\times(3\times10^{-2})=7{,}2\times10^{3}$ ;\
$(6\times10^{8})\div(1{,}5\times10^{3})=4\times10^{5}$.

\corrige{de l'exercice 9}
$3{,}7\times2{,}9=10{,}73$, donc $3{,}7\times10^{5}\times2{,}9\times10^{-8}=10{,}73\times10^{-3}=1{,}073\times10^{-2}$ :
ordre de grandeur $10^{-2}$.

\corrige{de l'exercice 10}
Estimation : $600\times5=3\,000=3\times10^{3}$, ordre de grandeur $10^{3}$. Valeur exacte :
$612\times4{,}8=2\,937{,}6$, bien de l'ordre de $10^{3}$.

\corrige{de l'exercice 11}
$2^{10}=1\,024$, donc $1\,000<2^{10}<10\,000$, soit $10^{3}<2^{10}<10^{4}$. Comme $1\,024$
est proche de $1\,000$, l'ordre de grandeur de $2^{10}$ est $10^{3}$.

\corrige{de l'exercice 12}
$2^{20}=\left(2^{10}\right)^{2}\approx\left(10^{3}\right)^{2}=10^{6}$. Valeur exacte :
$2^{20}=1\,048\,576$, très proche de $10^{6}$ (ordre de grandeur $10^{6}$).

\corrige{de l'exercice 13}
$149\,600\,000\text{ km}=1{,}496\times10^{8}$ km. Comme $1$ km $=10^{3}$ m, cette distance
vaut $1{,}496\times10^{8}\times10^{3}=1{,}496\times10^{11}$ m.

\corrige{de l'exercice 14}
$0{,}000007\text{ m}=7\times10^{-6}$ m. Nombre de globules pour $1$ cm $=10^{-2}$ m :
$\dfrac{10^{-2}}{7\times10^{-6}}\approx1\,429$ globules rouges.

\corrige{de l'exercice 15}
$\dfrac{2{,}7\times10^{7}}{8\times10^{9}}=0{,}3375\times10^{-2}=3{,}375\times10^{-3}$ :
ordre de grandeur $10^{-3}$ (la population du Cameroun représente environ un millième de la
population mondiale).

\corrige{de l'exercice 16}
$5{,}97\times10^{24}\text{ kg}=5{,}97\times10^{24}\div10^{3}\text{ t}=5{,}97\times10^{21}$
tonnes.

\corrige{du sujet 1}
1) $D=2^{3-5-(-4)}=2^{3-5+4}=2^{2}=4$.\\
2) $149\,600\,000=1{,}496\times10^{8}$ km.\\
3) $t=\dfrac{1{,}496\times10^{8}}{3\times10^{5}}\approx498{,}7$ s. En minutes :
$498{,}7\div60\approx8{,}31$ min, soit environ $8$ min $19$ s.\\
4) $E=\left(4\times10^{3}\right)^{2}=16\times10^{6}=1{,}6\times10^{7}$.

\corrige{du sujet 2}
1) $2^{20}=\left(2^{10}\right)^{2}=1\,024^{2}=1\,048\,576$ ;\
$2^{30}=2^{20}\times2^{10}=1\,048\,576\times1\,024=1\,073\,741\,824$.\\
2) Nombre de photos $=2^{30}\div2^{22}=2^{8}=256$.\\
3) $2^{30}=1\,073\,741\,824\approx1{,}07\times10^{9}$.\\
4) Ordre de grandeur : $10^{9}$ octets.

\corrige{du sujet 3}
1) $F=10^{7-2-(-3)}=10^{7-2+3}=10^{8}$.\\
2a) $\dfrac{2{,}7\times10^{7}}{8\times10^{9}}=0{,}3375\times10^{-2}=3{,}375\times10^{-3}\approx3{,}4\times10^{-3}$.\\
2b) En pourcentage : $3{,}375\times10^{-3}\times100=0{,}3375\,\%\approx0{,}3\,\%$.\\
3) $475\,000\text{ km}^{2}=4{,}75\times10^{5}\text{ km}^{2}$.

\end{multicols}
\exercicesBanner
